% Introduction: EMA with the idea of a moving average, MACD/SIGNAL, buy and sell
% signals.

\begin{frame}
  \frametitle{EMA to średnia krocząca, która najmocniej waży ostatnie sesje}
  EMA wygładza ceny. Uwzględnia wszystkie wcześniejsze sesje, ale ich wagi
  maleją wykładniczo, więc najnowsza cena znaczy najwięcej.
  \begin{equation*}
    \alpha_N = \frac{2}{N+1},\qquad
    E_0^{(N)} = x_0,\qquad
    E_i^{(N)} = \alpha_N \cdot x_i + (1 - \alpha_N) \cdot E_{i-1}^{(N)}
    \quad\text{dla } i \ge 1
    \tag{1}
  \end{equation*}
  \vskip2mm
  \begin{itemize}
    \item $x_i$ - cena zamknięcia z sesji numer $i$
    \item $E_i^{(N)}$ - wartość EMA o okresie $N$ po sesji $i$
    \item $\alpha_N$ - waga najnowszej ceny, tym większa, im krótszy
      okres $N$ (krótki okres daje szybką EMA, długi wolną)
  \end{itemize}
\end{frame}

\begin{frame}
  \frametitle{MACD to różnica szybkiej i wolnej EMA, a SIGNAL to wygładzone
    MACD}
  MACD to wskaźnik analizy technicznej, który na podstawie cen zamknięcia
  pomaga ocenić kierunek i dynamikę zmian cen oraz wyznacza sygnały kupna
  i sprzedaży.
  \begin{equation*}
    M_i = E_i^{(N_1)} - E_i^{(N_2)}
    \tag{2}
  \end{equation*}
  \begin{equation*}
    S_0 = M_0,\qquad
    S_i = \alpha_{N_3} \cdot M_i + (1 - \alpha_{N_3}) \cdot S_{i-1}
    \quad\text{dla } i \ge 1
    \tag{3}
  \end{equation*}
  \vskip2mm
  \begin{itemize}
    \item $M$ - linia MACD, różnica szybkiej EMA (okres $N_1$) i wolnej
      EMA (okres $N_2$)
    \item $S$ - linia SIGNAL, czyli EMA liczona z wartości MACD
      (okres $N_3$; wzór (1) dla $N = N_3$)
    \item Standardowe ustawienie: MACD(12, 26, 9).
  \end{itemize}
\end{frame}

\begin{frame}
  \frametitle{Zmiana znaku histogramu wyznacza sygnały kupna i sprzedaży}
  \begin{equation*}
    H_i = M_i - S_i
    \tag{4}
  \end{equation*}
  \vskip2mm
  \begin{itemize}
    \item Sygnał kupna w sesji $i$, gdy MACD przecina SIGNAL od dołu, czyli
      histogram przechodzi z wartości ujemnej lub zera na dodatnią
      ($H_{i-1} \le 0$ i $H_i > 0$).
    \item Sygnał sprzedaży w sesji $i$, gdy MACD przecina SIGNAL od góry,
      czyli histogram przechodzi z wartości dodatniej lub zera na ujemną
      ($H_{i-1} \ge 0$ i $H_i < 0$).
    \item Transakcja zapada w sesji po sygnale, kupuje się za cały kapitał
      i sprzedaje całą pozycję. Kapitał początkowy to
      $C_{100} = 1000 \cdot x_{100}$, a sygnały liczy się od próbki
      \resFirstSignalIndex{}.
  \end{itemize}
\end{frame}
